% !TEX root = ../monograph.tex
\subsection{Constructing the restarted representations}
\label{ch11:restart-construction}

We now construct the auxiliary history just described. The physical
insertions retain their original calibration; only the integration
below the comparison level is restarted. All support norms in this
subsection use connected covers and include full dependence regions.
Their exponents are chosen with the reserves required by the marked
expansion before choosing an inverse coupling threshold.

Write $x_i=\widehat\beta_i$ and, for $x$ large, put
\[
 H(x)=e^{(\log\log x)^2},\quad h(x)=e^{(\log\log x)^{3/2}},\quad
 b(x)=\lceil(\log x)^2\rceil,\quad
 \vartheta(x)=x^{-1/2}h(x).
\]
The new representation at level $i$ has radius $B_i=b(x_i)^2$
and threshold $\vartheta(x_i)$. Its clip has the original two
coefficients, evaluated at this new threshold. Its path coefficients
and block-kernel branches are the original ones. In particular,
``retaining the clip'' means retaining its coefficients and formula,
not asserting equality of two different threshold functions.

\paragraph{Choice of reference parameter for the cutoff.}
There is a unique sufficiently large $r(x)$ satisfying
\[
 r(x)^{-1/2}H(r(x))=\vartheta(x).
\]
Indeed the function on the left is eventually strictly decreasing;
its values at $x$ and $x^2$ bracket the right side. Thus
$x\le r(x)\le x^2$, and taking logarithms gives
$r(x)=x^{1+o(1)}$. With this choice the original threshold formula evaluated at $r(x)$
gives the new threshold at $x$. The inverse coupling in the block
integral still lies in the band determined by $x$. In particular
\[
 b(x)\le b(r(x))\le4b(x),\qquad
 25b(r(x))+30\le B_i\le b(r(x))^2
\]
eventually. On each fixed adjacent stage pair,
$x_{i+1}=x_i+\kappa_G\log L_i$ implies
$B_{i+1}\le2B_i$ and
$\vartheta(x_{i+1})/\vartheta(x_i)\to1$.
The local chart, clipped source, Gaussian precision and tail events
are those of Chapter~\ref{ch06:large-fields} with reference parameter $r(x_i)$
and threshold $\vartheta(x_i)$. Its elementary scale inequalities
continue to hold: the radii are logarithmic powers and the thresholds
are $x^{-1/2+o(1)}$. The actual inverse coupling band remains expressed in
$x_i$, not in $r(x_i)$. In particular one must not invoke a lower
bound $\beta\ge r(x_i)/2$. For every fixed $\eta>0$, changing the reference parameter and
the actual band instead give, eventually,
\[
 y^{1-\eta}\le\beta\le2y,\qquad y=r(x_i).
\]
For a prescribed finite local estimate, fix the number of powers
of $\beta^{-1}$ and $g=\beta^{-1/2}$ that occur before choosing
$\eta$. Replacing a comparison power of $y$ by the corresponding
actual inverse coupling power then loses at most $y^{C\eta}$, where $C$
is fixed by that estimate. Choose $\eta$ small enough to retain
its negative exponent. The dimensionless quadratic gap and path
bounds are unchanged; physical covariance factors retain their
actual $\beta^{-1}$ normalization. Likewise, the cutoff-to-noise
ratio is computed exactly as
$\vartheta(x_i)/g\asymp h(x_i)$, and the selected-component energy
is $\beta\vartheta(x_{i+1})^2\ge c h(x_{i+1})^2$.
These comparisons, not a fictitious inverse coupling band in $y$, are
used below.

\begin{proposition}[Re-expression of the endpoint]
\label{ch11:restart-endpoint}
At a sufficiently late comparison level $k$, the original density has an
exact presentation with radius $b(x_k)^2$, threshold
$\vartheta(x_k)$, the original path coefficients, clip coefficients,
inverse coupling and scalar, regular bound no larger than its original
bound, and large-field site bound
\[
 G(x_k)=\exp\bigl(-a_*h(x_k)^2\bigr).
\]
Here $a_*>0$ is chosen sufficiently small at each fixed stage before
its threshold for the starting scale. The construction retains the old large-field
labels and their eligibility criteria inside the new activities.
\end{proposition}
\begin{proof}
Abbreviate $s=x_k$, $t=s^{-1/2}H(s)$, $\vartheta=\vartheta(s)$,
$M=b(s)$ and $B=M^2$. Write $A_t,A_\vartheta$ for the path actions
with their respective event indicators and clips, and $E_t,E_\vartheta$ for
the original regular coefficients with the corresponding event indicators.
The old regular bound is at most $s^{-11/10}$; the old
large-field bound is at most $w_s=e^{-H(s)^{1/5}}$.
These are the selected trajectory's endpoint bounds.

Eventually $2\vartheta\le t$. The clip comparison proved above gives
$d_t(r)-d_\vartheta(r)\ge c_{\rm lin}\vartheta^2$ whenever
$r>\vartheta$, and a nonnegative difference otherwise.
Distribute each plaquette's difference among its incident sites.
Allowing five sites is a harmless upper bound. For the set
$F_\vartheta$ of new large-field sites this gives local energy contributions $d_z$ with
\[
 \sum_zd_z=\sum_p(d_t(r_p)-d_\vartheta(r_p)),\qquad
 d_z\ge c_{\rm lin}\vartheta^2/5\quad(z\in F_\vartheta).
\]
The omitted quadratic path rows give another nonnegative energy contribution.
Truncate each row at true path length $B$ before squaring, and keep
its half-square only where the old retention condition holds and the new one
fails. The relaxed small-field condition reads a radius-$B$ neighborhood: outside it
$2\le\vartheta e^{B/c}$ makes the condition automatic. The truncated
half-square therefore reads radius at most $2B$ and vanishes if its
tested neighborhood is new-small-field. Distribute its value equally
among the new large-field sites in that neighborhood, making no division
when the set is empty. The resulting local energy contributions $b'_z\ge0$ depend only on a neighborhood of radius
at most $4B$, and their sum is exactly the sum of truncated
half-squares. Perform this operation separately for each signed
orientation, retaining its averaging weight.

Choose $a_*\le c_{\rm lin}/20$, put $u=a_*h(s)^2$, and for large-field sites
$z$ put $X_z=\exp[-\beta(d_z+b'_z)]$. The actual band gives
$\beta\vartheta^2\ge h(s)^2/2$, hence $|X_z|\le e^{-2u}$.
The remaining exponent difference is
\[
 D=-(A_t-A_\vartheta)+\beta\sum_z(d_z+b'_z)
                               +(E_\vartheta-E_t).
\]
It consists of the omitted quadratic row pairs, the cubic, quartic
and loop event indicator differences, and the regular event indicator differences.
Each term vanishes on a new-small-field full dependence region.
At a fixed larger support exponent $a+2$ its site bound is
\begin{equation}
 \epsilon_s\le C B^C\bigl(
 \beta C_{\rm path}^2e^{-\sigma B/4}
 +\Gamma_s t^2+\beta C_{\rm path}(t^3+t^4)+s^{-11/10}\bigr)
                  =s^{-1/2+o(1)}.
 \label{ch11:restart-endpoint-error}
\end{equation}
Here $\sigma>0$ is a retained true-length path rate and
$\Gamma_s\le2H(s)$ is the loop coefficient bound.
For the first term expand the difference of squares into ordered
path pairs, one of length exceeding $B$, and retain part of their
exponential length weight. For the atoms restricted by event indicators use their old
bounds $t^j$ and the cover bound $1+2j+\operatorname{length}/B$.
For an old support $Y$, its fixed halo obeys
$c_B(Y^{[8]})\le C+C(M/B)c_M(Y)$. Eventually its old exponent
absorbs this new exponent and the extra unit used to turn an anchored
bound into a site bound. These are the support conversions of
Section~\ref{ch06:primitive-bounds}; their losses are polynomial in
$B$. The same argument bounds enlarged old large-field labels by
$CB^Cw_s$. No path multiplicity or winding path is removed.

Coalesce the absolutely summable terms by full support into $D_Y$.
For large $s$, put $W_Y=e^{D_Y}-1$, so $|W_Y|\le2|D_Y|$ and
$W_Y$ vanishes when its support is new-small-field. The identity
\begin{equation}
 e^{-A_t-E_t}=e^{-A_\vartheta-E_\vartheta}
                \prod_Y(1+W_Y)\prod_{z\in F_\vartheta}X_z
 \label{ch11:restart-endpoint-identity}
\end{equation}
is exact. Multiply by the old large-field gas and expand the optional
$W_Y$ factors. Keep the old compatibility and eligibility criteria.
Group using intersections of full dependence supports, including the fixed
halo for old large-field conditions and the radius-$4B$ supports of $X_z$.
Every nonzero component contains a selected $X_z$: an old activity
contains an old large-field site, hence a new large-field site, and every nonzero $W_Y$
reads a new large-field indicator. The factor at that site is mandatory and intersects
that component. Assign the component its exact new set of large-field sites and also
retain the union of its old sets of large-field sites and all their conditions.
Conversely, those constituents recover the original compatible
collection and optional subset. Thus the regrouping is a bijection
of terms, not a replacement of old eligibility by a new condition.
Disjoint new exclusion supports are the unions of these full read
supports. Their eligibility criteria require exact coverage of the new large-field sites there.

Reserve $e^{-u}$ per selected site factor. The remaining site sum
at exponent $a+1$ is at most
\[
 \rho_s=2\epsilon_s+CB^Cw_s+
       (8B+1)^4e^{(a+1)9^4}e^{-u}=s^{-1/2+o(1)}.
\]
A radius-$4B$ box has a connected $B$-cover of at most $9^4$
boxes; at most $(8B+1)^4$ such boxes contain a given site.
The connected-set bound of Chapter~\ref{ch08:chapter}, losing one
support exponent, applies since $(2B+1)^4\rho_s\le1$.
Every component retains at least one reserved factor, so its total
site bound is at most $e^{-u}(e^{\rho_s}-1)\le G(s)$.
There is no extra sum over prescribed sets of large-field sites: selected site factors already
specify them, and exact-coverage conditions only remove terms.
Absolute convergence justifies countable labels and coalescing.

The regular coefficients themselves are unchanged. Their domains
and allowed directions shrink with the threshold, and increasing
the cover radius decreases their support weights. Their old bounds
therefore persist. The exact identity and the gas bijection prove
the assertion, with unchanged scalar. All operations preserve gauge
invariance, reality and signed symmetry, by reindexing the original
orientation average.
\end{proof}

\paragraph{The genuine density and its neutral reference.}
At an intermediate level write the genuine presentation as
$R=e^{-E}R^0$, where $R^0$ is obtained by deleting only the regular
additive family $E=\sum_\alpha E_\alpha$. Its path, clip, large-field
activities and all their conditions remain unchanged. If the anchored
weighted value bound of that family is $q$, then
$\sum_\alpha|E_\alpha|\le q|\Lambda|$; this proves integrability
at finite volume, but is not used as a volume-uniform estimate.
Instead insert local factors
\[
 h_\alpha=e^{-E_\alpha}-1,\qquad
 \prod_\alpha(1+h_\alpha)=e^{-E}.
\]
Since $|h_\alpha|\le e^{|E_\alpha|}|E_\alpha|$, their local site
bound is the old anchored bound times a fixed polynomial in the
radius. They are inserted into the same expansion as the neutral
Gaussian integral. Thus the object integrated is always $R$, also
for every coefficient of a marked density. Deleting $E$ for the
neutral estimates never deletes it from the density identity.

\begin{proposition}[One restarted history on the comparison window]
\label{ch11:restart-family}
Fix the finite degree and derivative requests and the comparison
window with top level $k$, reference inverse coupling $s=x_k$, and at most
$(\log\log s)^4$ steps. Take its thresholds for the starting scales large enough
that $3s/4\le x_i\le s$ and
$|\beta_{N,p,k}-s|\le s/8$. On the original adjacent path data there
are presentations $P_i$, $i\le k$ in that window, with the bounds
$(B_i,\vartheta(x_i))$, common actual inverse coupling
$\beta=\beta_{N,p,k}$, large-field bound $G(x_i)$ and additive
value cost at most
\[
 q_{k-i}=s^{-1/200}A_s^{4(k-i)},\qquad
                    A_s=e^{(\log\log s)^3}.
\]
Their exact densities satisfy
\begin{equation}
 D_k=\rho_{N,p,k},\qquad
 D_i=\mathcal B_{i,\beta}D_{i+1}=\operatorname{density}(P_i),
 \qquad \int D_i=Z_{N,p}.
 \label{ch11:restart-density}
\end{equation}
Here the block operator has the original geometric branches, evaluated
at the common inverse coupling. The same local factors, normalization and
all deletion-of-support identities supply both weighted marked-step
uses. The constants are uniform in cutoff, period and the compatible
local fields, for the fixed requests.
\end{proposition}
\begin{proof}
Start with Proposition~\ref{ch11:restart-endpoint}.
The reference-inverse-coupling band gives $x_i/2\le\beta\le2x_i$ on every step.
Only the actual inverse coupling is frozen; the reference inverse couplings, original path pairs,
metric floors, gauge-covariant block average derivatives and forest charts keep their
original adjacent values. The recording-scale inequalities above
verify the radii and threshold comparisons in the local integral.
The fine path and clip in that integral are exactly the ones in the
incoming presentation. A change of stage uses its own adjacent pair,
not the previous stage's quadratic form.

Suppose the genuine incoming value cost is $q_n$, $n=k-i-1$.
Its neutral version has zero such cost and the same large-field
bound and loop coefficient bound. Thus it meets the small incoming
regular bound in Chapter~\ref{ch06:large-fields}; we do not claim
that the larger $q_n$ meets that derivative bound. The genuine
additive factors are handled by the local multipliers just defined.

We spell out the quantitative margins on this actual expansion.
The primitive exponent decomposition in
Section~\ref{ch06:primitive-bounds}, with reference parameter
$y=r(x_i)$, gives an unselected local exponent of site size
$t_y$ with $2t_y\le y^{-1/48}$ after enlargement of its threshold.
Indeed the corresponding exterior row, determinant and Haar
bounds are $y^{-1/4+o(1)}$, including radius polynomials and the
actual inverse coupling factors just discussed; the linear residual is
bounded by $y^{-4}y^2$ before these subpower losses. The two
determinant-word series and inverse correction are exponentially
small in $b(y)$, including the $y^4$ factor for two Gaussian
coordinates and any fixed inverse powers of the actual inverse coupling. The inherited path
terms have costs $O(\Gamma\vartheta^2)$ and
$O(\beta(\vartheta^3+\vartheta^4))$, with true-length reserves.
The inherited bounds before changing the reference parameter can be weakened to
$y^{-1/12}$, with fixed finite powers of the actual inverse coupling;
choose the preceding $\eta$ so their loss is less than $1/48$.
Thus all these are negative powers with a strict margin beyond
$1/24$. One may weaken their sum successively to $y^{-1/24}$
and then $y^{-1/48}$
to include exponentiation, support conversion and fixed constants.
The neutral additive part is zero. The local exponents are bounded
extensions on the same chart event indicators; outside them the failed cutoff conditions
are retained as tail-event factors, not evaluated as small exponents.

Selected components have, after the extracted Gaussian and local
normalization factors have been bounded, bound
$\exp[-K h(x_{i+1})^2 r]$ for $r$ selected cutoff conditions, by the actual
component energy in Chapter~\ref{ch06:large-fields}. Choose
$a_*>0$ also so small that $4a_*\le K/2$.
For both marked-step uses at once take
\[
 m=\max\{G(x_i)^{-2},G(x_{i+1})^{-1}\},\qquad
                 1\le l\le\max(1,H(s)^{q-3}),
\]
where $q$ is the fixed requested finite power.
The ratio $h(x_i)^2/h(x_{i+1})^2$ tends to one uniformly in the
window: its logarithm is $O(\sqrt{\log\log s}/\log s)$.
Thus $\log m\le4a_*h(x_{i+1})^2$ eventually, retaining at least
half the component energy. Apply the factor $m$ only to components with an
actual selected prescribed set of large-field sites, and to a selected component within each
such component. The original prescribed set of large-field sites coverage ensures that an
incoming large-field factor cannot occur without such a component.
Factors $l$, finite decorations and local counting costs have
logarithms $O((\log\log s)^2+\log B_i)$ per selected cutoff condition and fit
inside the remaining margin.

The simultaneous tail estimate of Chapter~\ref{ch06:large-fields}
uses these very chart coordinates and event indicators under the same Gaussian
law, including any bounded genuine multipliers. It therefore bounds
all compatible factors from failed cutoff conditions jointly. Its cutoff exponent
and the retained selected-component exponent dominate every multiple
of $\log s$. The unselected components after inflation cost at most
$s^{-1/48+o(1)}$. The projected old large-field site sum is bounded
by ${\cal W}(B_i)L_i^4T_eG(x_{i+1})$, also smaller than every inverse
power of $s$. It need not include $m$ a second time: the selected
component estimate already includes that factor.

The genuine additive multipliers have uninflated site bound
$C(42B_i+1)^4q_n$. Every fixed product of radius factors, support
conversion constants and $l$ is eventually at most $A_s^{1/4}$.
Hence their inflated site bound is at most $q_nA_s^{1/4}$.
For $n\le(\log\log s)^4$ and every fixed $r$,
\[
 q_nA_s^r\le s^{-1/400}.
\]
This follows by taking logarithms: the positive contribution is
$O((\log\log s)^7)$, whereas the negative term is
$-(\log s)/200$. The connected expansion can therefore be rooted
at an additive multiplier when one occurs; its site bound fits
$\omega_n=q_nA_s$. The other components fit
$\delta_n-\omega_n$, where $\delta_n=q_nA_s^3$, since
$s^{-1/48+o(1)}=o(q_n)$. This proves the inflated product bound and
its site summability on the genuine, not a substituted, density.
All finite stage factors in this paragraph are absorbed simultaneously
by its choice of a sufficiently large starting scale; no uniform supremum over infinitely many stages
is used.

It remains to reconstruct the coarse presentation. The regular
connected logarithm has anchored cost at most $16\delta_n$.
Keep the original coarse path action $A_i$. Its difference from
the base action extracted in the Gaussian integral consists of the
cubic, quartic and quadratic-loop terms restricted by their event indicators, with cost at most
\[
 \beta e^{7e}C_i\vartheta(x_i)^3+
 \beta e^{9e}C_i\vartheta(x_i)^4+
 e^{5e}\Gamma_i\vartheta(x_i)^2\le s^{-1/200}.
\]
The cover bound for a degree-$j$ path atom is
$1+2j+\operatorname{length}/B_i$; the path decay absorbs the last term.
Add this correction to the regular logarithm. Its total value cost
is at most $s^{-1/200}+16\delta_n\le q_{n+1}$.
The sign is fixed by
$e^{-\beta B_i^{\rm ref}}e^{-\ell_i}
=e^{-A_i}e^{-(\ell_i+\beta B_i^{\rm ref}-A_i)}$;
$B_i^{\rm ref}$ here is the base action, not the radius.

At the sliced-integrand level the event indicators retain the finite tag-count
factor; it cancels the denominator of the normalized tag average
when the slices are summed. Signed orientations retain their original
probability weights. Neither averaging introduces a second determinant
or a second extracted scalar.
The positive Gaussian prefactor determines a real scalar logarithm.
No branch of a complex logarithm is chosen. The remaining
large-field factors, including the regular logarithms with holes,
are grouped at their full support and exact old prescribed sets of large-field sites as in
\eqref{ch08:reassembly}. The raised hole correction is bounded by
$16l\delta_n e^{16\delta_n}$, using
$|e^{-c}-1|\le e^{|c|}|c|$.
This and the other component costs are $s^{-1/400+o(1)}$, within
the available $H(s)^{-1/8}$ reserve. The inflated connected bound
gives outgoing large-field bound at most $m^{-1}\le G(x_i)^2$,
hence at most $G(x_i)$. No residual path correction is left outside
the regular family. Averaging the regular pieces over the finite
signed motion group retains their anchored bound and invariant sum.
The exact factorization and reassembly thus give
$D_i=\mathcal B_{i,\beta}D_{i+1}$ on every coarse configuration,
including configurations with failed cutoff conditions. There is no division
by a large-field activity, which can vanish.

At each step choose this expansion before specifying the marked
family and retain it for subsequent use. Smaller support exponents
and smaller inflations are then estimates on the same coefficients.
Only the newly constructed lower entry is installed; the endpoint
and all already installed higher entries are unchanged. This proves
finite descending induction, including adjacent stage changes.
Finally positivity of $D_i$ follows from the positive endpoint and
probability kernels, not from positivity of individual activities.
Their mass-one identity and finite-volume integrability give
$\int D_i=Z_{N,p}$ by Fubini. The same equalities hold coefficientwise
for finite marked polynomials. This proves the asserted normalization
and the common-history statement.
\end{proof}
